\vspace{0.5em}\noindent\textbf{Chia sẻ và phản
hồi}\par

Phần này tổng hợp những chia sẻ cá nhân của mình sau khi tham gia chương
trình \textbf{First Cloud AI Journey (FCAJ)} và thực hiện project
\textbf{Internship Application Platform}. Nội dung tập trung vào trải
nghiệm học tập, sự hỗ trợ từ mentor/cộng đồng, những điểm mình đánh giá
cao và các đề xuất để chương trình có thể cải thiện hơn trong tương lai.

\vspace{0.5em}\noindent\textbf{Đánh giá chung}\par

\vspace{0.5em}\noindent\textbf{1. Môi trường học tập và thực
hành}\par

Môi trường của chương trình tạo điều kiện tốt để mình tiếp cận các chủ
đề Cloud và AI theo hướng thực tế. Thay vì chỉ học lý thuyết, mình có cơ
hội xây dựng một project có nhiều thành phần như frontend, backend, chat
realtime, AI service, database, Docker, Kubernetes, CI/CD và
observability.

Điều mình đánh giá cao là chương trình khuyến khích người học tự tìm
hiểu, tự triển khai và tự ghi nhận kết quả. Cách tiếp cận này giúp mình
hiểu sâu hơn về cách một hệ thống được thiết kế, vận hành và triển khai
trên môi trường Cloud.

\vspace{0.5em}\noindent\textbf{2. Sự hỗ trợ từ mentor và team
admin}\par

Mentor và team admin hỗ trợ tốt trong quá trình học tập. Khi gặp vấn đề,
mình có thể đặt câu hỏi, nhận định hướng và tiếp tục tự xử lý thay vì
chỉ nhận đáp án có sẵn. Điều này giúp mình rèn luyện tư duy debug, đọc
tài liệu và giải quyết vấn đề độc lập.

Các thông tin về chương trình, sự kiện, workshop và tài liệu tham khảo
cũng được chia sẻ tương đối rõ ràng, giúp mình có thêm nguồn học phù hợp
cho từng giai đoạn.

\vspace{0.5em}\noindent\textbf{3. Sự phù hợp giữa công việc và chuyên ngành
học}\par

Nội dung thực tập phù hợp với chuyên ngành Computer Science và định
hướng phát triển của mình. Project có nhiều phần liên quan đến backend,
database, distributed systems, Cloud infrastructure, security và AI
integration.

Thông qua project, mình không chỉ củng cố kiến thức lập trình mà còn
hiểu thêm về các yếu tố phi chức năng như bảo mật, khả năng mở rộng, độ
tin cậy, monitoring và kiểm soát chi phí.

\vspace{0.5em}\noindent\textbf{4. Cơ hội học hỏi và phát triển kỹ
năng}\par

Trong quá trình thực tập, mình học được nhiều kỹ năng quan trọng:

\begin{itemize}
\tightlist
\item
  Thiết kế hệ thống nhiều dịch vụ.
\item
  Xây dựng REST API với FastAPI.
\item
  Quản lý dữ liệu bằng PostgreSQL, DynamoDB, Redis và S3.
\item
  Tích hợp realtime chat bằng Socket.IO.
\item
  Tích hợp AI để phân tích CV và Job Description.
\item
  Container hóa bằng Docker và triển khai bằng Kubernetes.
\item
  Thiết lập CI/CD và quan sát hệ thống bằng metrics, logs, traces.
\item
  Viết báo cáo kỹ thuật song ngữ theo cấu trúc rõ ràng.
\end{itemize}

Những kỹ năng này giúp mình có cái nhìn thực tế hơn về cách xây dựng một
hệ thống Cloud-native thay vì chỉ tập trung vào từng chức năng riêng lẻ.

\vspace{0.5em}\noindent\textbf{5. Văn hóa học tập và tinh thần cộng
đồng}\par

Cộng đồng FCAJ tạo cảm giác cởi mở và khuyến khích chia sẻ. Các sự kiện
trực tuyến, workshop và bài blog trong nhóm giúp mình có thêm góc nhìn
từ nhiều chủ đề khác nhau như AWS Cloud, Agentic AI, kiến trúc doanh
nghiệp, DevOps, observability và bảo mật.

Mình cũng đánh giá cao việc chương trình khuyến khích người học chia sẻ
lại kiến thức thông qua blog và báo cáo. Đây là cách tốt để biến kiến
thức đã học thành nội dung có thể trình bày và kiểm chứng.

\vspace{0.5em}\noindent\textbf{Một số câu hỏi
khác}\par

\vspace{0.5em}\noindent\textbf{Điều mình hài lòng nhất trong thời gian thực tập là
gì?}\par

Điều mình hài lòng nhất là có cơ hội xây dựng một project tương đối đầy
đủ, bao gồm cả business flow, AI integration, realtime communication,
deployment direction và observability. Project giúp mình hiểu rõ hơn mối
liên hệ giữa code, hạ tầng và vận hành.

\vspace{0.5em}\noindent\textbf{Điều mình nghĩ chương trình có thể cải thiện cho các thực
tập sinh sau là
gì?}\par

Chương trình có thể bổ sung thêm checklist theo từng giai đoạn, ví dụ
checklist cho backend, frontend, AWS deployment, monitoring và cost
cleanup. Điều này sẽ giúp người học dễ tự kiểm tra tiến độ và tránh bỏ
sót những phần quan trọng.

Ngoài ra, nếu có thêm một số buổi review project mẫu hoặc phân tích lỗi
thường gặp khi triển khai trên AWS, người học sẽ có thêm kinh nghiệm
thực tế trước khi tự làm project.

\vspace{0.5em}\noindent\textbf{Nếu giới thiệu cho bạn bè, mình có khuyên họ tham gia
chương trình
không?}\par

Mình sẽ khuyên các bạn quan tâm đến Cloud, backend, DevOps hoặc AI
application tham gia chương trình. FCAJ phù hợp với những người muốn học
theo hướng thực hành, tự xây dựng sản phẩm và hiểu cách các dịch vụ AWS
có thể kết hợp trong một hệ thống thực tế.

\vspace{0.5em}\noindent\textbf{Đề xuất và mong
muốn}\par

\begin{itemize}
\tightlist
\item
  Có thêm project checklist theo từng milestone để người học dễ theo dõi
  tiến độ.
\item
  Có thêm ví dụ về cost optimization và cleanup để tránh phát sinh chi
  phí không cần thiết trên AWS.
\item
  Có thêm buổi chia sẻ về lỗi phổ biến khi dùng EKS, IAM, S3,
  CloudFront, RDS và GitHub Actions.
\item
  Có thêm guideline về cách viết báo cáo kỹ thuật và cách trình bày
  architecture diagram.
\item
  Có thể tổ chức thêm phiên project review ngắn để người học nhận góp ý
  trước khi hoàn thiện báo cáo.
\item
  Có thêm định hướng về career path cho Cloud Engineer, Backend
  Engineer, DevOps Engineer và AI Engineer.
\end{itemize}

\vspace{0.5em}\noindent\textbf{Tác động đến định hướng nghề
nghiệp}\par

Sau chương trình, mình có định hướng rõ hơn về việc theo đuổi mảng
Cloud/Backend Engineering kết hợp AI integration. Mình nhận ra rằng việc
xây dựng một ứng dụng hiện đại không chỉ là viết code, mà còn bao gồm
thiết kế kiến trúc, bảo mật, triển khai, monitoring, quản lý chi phí và
chuẩn bị tài liệu vận hành.

Chương trình giúp mình tự tin hơn khi đọc tài liệu kỹ thuật, phân tích
yêu cầu, lựa chọn dịch vụ phù hợp và trình bày giải pháp theo hướng có
cấu trúc.

\vspace{0.5em}\noindent\textbf{Lời nhắn cuối}\par

Mình cảm ơn chương trình First Cloud AI Journey, mentor, team admin và
cộng đồng đã tạo điều kiện để mình học hỏi và hoàn thiện project trong
kỳ thực tập. Những kiến thức và trải nghiệm này là nền tảng quan trọng
để mình tiếp tục phát triển kỹ năng Cloud, backend, DevOps và AI
application trong thời gian tới.
